\begin{longtable}{|p{0.8cm}|p{7.5cm}|>{\centering\arraybackslash}p{0.8cm}|>{\centering\arraybackslash}p{0.8cm}|>{\centering\arraybackslash}p{0.8cm}|>{\centering\arraybackslash}p{0.8cm}|>{\centering\arraybackslash}p{0.8cm}|}
    \caption{Hasil Kuesioner \textit{User Acceptance Test} - Perspektif \textit{User}} \label{tabel:uat_user_hasil}                                                                                                                                   \\
    \hline
    \multirow{2}{*}{\textbf{No.}}              & \multirow{2}{*}{\textbf{Pertanyaan}}                                                                   & \multicolumn{5}{c|}{\textbf{Jawaban}}                                                       \\ \cline{3-7}
                                               &                                                                                                        & \textbf{STS}                          & \textbf{TS} & \textbf{N} & \textbf{S} & \textbf{SS} \\
    \hline
    \endfirsthead

    \multicolumn{7}{c}%
    {{\bfseries Tabel \thetable\ lanjut dari halaman sebelumnya}}                                                                                                                                                                                     \\
    \hline
    \multirow{2}{*}{\textbf{No.}}              & \multirow{2}{*}{\textbf{Pertanyaan}}                                                                   & \multicolumn{5}{c|}{\textbf{Jawaban}}                                                       \\ \cline{3-7}
                                               &                                                                                                        & \textbf{STS}                          & \textbf{TS} & \textbf{N} & \textbf{S} & \textbf{SS} \\
    \hline
    \endhead

    \hline \multicolumn{7}{r}{\textit{Bersambung ke halaman berikutnya}}                                                                                                                                                                              \\
    \endfoot

    \hline % \endlastfoot akan menjadi baris terakhir tabel setelah "Total Skor"
    \endlastfoot

    % --- Pertanyaan UAT untuk Pengguna (10 Pertanyaan) ---
    1.                                         & Apakah proses pendaftaran akun dan login ke aplikasi Tamu.id mudah dipahami dan dilakukan?             &                                       &             &            & 1          & 9           \\
    \hline
    2.                                         & Apakah tampilan antarmuka aplikasi Tamu.id menarik dan mudah digunakan secara umum?                    &                                       &             & 1          & 3          & 6           \\
    \hline
    3.                                         & Apakah informasi mengenai cabang dan tipe pod (harga, fasilitas, ketersediaan) disajikan dengan jelas? &                                       &             &            & 3          & 7           \\
    \hline
    4.                                         & Apakah alur proses pencarian dan pemesanan pod mudah diikuti dan efisien?                              &                                       &             &            & 3          & 7           \\
    \hline
    5.                                         & Apakah proses pembayaran saat melakukan pemesanan mudah dan terasa aman?                               &                                       &             & 1          & 1          & 8           \\
    \hline
    6.                                         & Apakah proses check-in dan check-out melalui aplikasi berjalan dengan lancar?                          &                                       &             &            &            & 10          \\
    \hline
    7.                                         & Apakah penggunaan QR Code untuk akses pod (jika ada fitur ini) mudah dan praktis?                      &                                       &             &            & 1          & 9           \\
    \hline
    8.                                         & Apakah saya mudah menemukan riwayat pemesanan dan detail transaksi saya?                               &                                       &             & 1          & 3          & 6           \\
    \hline
    9.                                         & Apakah fitur untuk memberikan ulasan dan rating setelah menginap mudah digunakan?                      &                                       &             &            & 5          & 5           \\
    \hline
    10.                                        & Secara keseluruhan, apakah Anda puas dengan fungsionalitas dan kemudahan penggunaan aplikasi Tamu.id?  &                                       &             &            & 1          & 9           \\
    \hline
    % --- Baris Total Bobot, Skor, Total Skor (Format Baru) ---
    \multicolumn{2}{|l|}{\textbf{Total Bobot}} &                                                                                                        &                                       & 3           & 21         & 76                       \\
    \hline
    \multicolumn{2}{|l|}{\textbf{Skor}}        &                                                                                                        &                                       & 9           & 84         & 380                      \\
    \hline
    \multicolumn{2}{|l|}{\textbf{Total Skor}}  & \multicolumn{5}{c|}{473}                                                                                                                                                                             \\ % Mengganti "anjay" dengan "..."
\end{longtable}

Dari tabel \ref{tabel:uat_user_hasil} di atas, kita dapat menghitung skor \textit{User Acceptance Test} (UAT) untuk perspektif pengguna sebagai berikut:

\begin{align*}
    \text{Skor UAT User}(\%) & = \left( \frac{A}{B \times N} \right) \times 100\%     \\
                             & = \left( \frac{473}{50 \times 10} \right) \times 100\% \\
                             & = \left( \frac{473}{500} \right) \times 100\%          \\
                             & = 0.946 \times 100\%                                   \\
                             & = 94.6\%
\end{align*}

Kesimpulan dari hasil UAT perspektif pengguna menunjukkan bahwa aplikasi Tamu.id mendapatkan skor 94.6\%, yang berarti aplikasi ini sangat diterima oleh pengguna dan masuk dalam kategori Sangat Layak berdasaarkan tabel \ref{tabel:kategori-kelayakan}.

\begin{longtable}{|p{0.8cm}|p{7.5cm}|>{\centering\arraybackslash}p{0.8cm}|>{\centering\arraybackslash}p{0.8cm}|>{\centering\arraybackslash}p{0.8cm}|>{\centering\arraybackslash}p{0.8cm}|>{\centering\arraybackslash}p{0.8cm}|}
    \caption{Hasil Kuesioner \textit{User Acceptance Test} - Perspektif \textit{Admin} Cabang} \label{tabel:uat_admin_cabang_hasil}                                                                                                                                                    \\
    \hline
    \multirow{2}{*}{\textbf{No.}}              & \multirow{2}{*}{\textbf{Pertanyaan}}                                                                                                    & \multicolumn{5}{c|}{\textbf{Jawaban}}                                                       \\ \cline{3-7}
                                               &                                                                                                                                         & \textbf{STS}                          & \textbf{TS} & \textbf{N} & \textbf{S} & \textbf{SS} \\
    \hline
    \endfirsthead

    \multicolumn{7}{c}%
    {{\bfseries Tabel \thetable\ lanjut dari halaman sebelumnya}}                                                                                                                                                                                                                      \\
    \hline
    \multirow{2}{*}{\textbf{No.}}              & \multirow{2}{*}{\textbf{Pertanyaan}}                                                                                                    & \multicolumn{5}{c|}{\textbf{Jawaban}}                                                       \\ \cline{3-7}
                                               &                                                                                                                                         & \textbf{STS}                          & \textbf{TS} & \textbf{N} & \textbf{S} & \textbf{SS} \\
    \hline
    \endhead

    \hline \multicolumn{7}{r}{\textit{Bersambung ke halaman berikutnya}}                                                                                                                                                                                                               \\
    \endfoot

    \hline % \endlastfoot
    \endlastfoot

    % --- Pertanyaan UAT untuk Admin Cabang (10 Pertanyaan) ---
    1.                                         & Apakah proses login ke dashboard Admin Cabang mudah dan aman?                                                                           &                                       &             &            &            & 1           \\
    \hline
    2.                                         & Apakah tampilan dashboard Admin Cabang informatif untuk memantau operasional harian cabang?                                             &                                       &             &            &            & 1           \\
    \hline
    3.                                         & Apakah pengelolaan informasi umum cabang (seperti detail kontak, foto, atau deskripsi) mudah dilakukan?                                 &                                       &             &            &            & 1           \\
    \hline
    4.                                         & Apakah fitur untuk menambah dan mengelola tipe pod serta unit pod individual di cabang saya mudah digunakan?                            &                                       &             &            &            & 1           \\
    \hline
    5.                                         & Apakah proses untuk mengatur harga pod (termasuk harga dinamis jika digunakan) cukup jelas dan mudah diterapkan?                        &                                       &             &            &            & 1           \\
    \hline
    6.                                         & Apakah pengelolaan data reservasi tamu (melihat daftar booking, detail, dan status) efisien?                                            &                                       &             &            &            & 1           \\
    \hline
    7.                                         & Apakah fitur untuk menginput atau mengelola data reservasi yang berasal dari Online Travel Agent (OTA) mudah dioperasikan?              &                                       &             &            &            & 1           \\
    \hline
    8.                                         & Apakah proses untuk mengelola status kebersihan dan ketersediaan pod (misalnya, dari "Perlu Dibersihkan" menjadi "Tersedia") sederhana? &                                       &             &            &            & 1           \\
    \hline
    9.                                         & Apakah fitur untuk melihat dan menanggapi (memoderasi) ulasan dari pengguna untuk cabang saya mudah digunakan?                          &                                       &             &            &            & 1           \\
    \hline
    10.                                        & Secara keseluruhan, apakah dashboard Admin Cabang ini efektif membantu saya dalam mengelola operasional cabang sehari-hari?             &                                       &             &            &            & 1           \\
    \hline
    % --- Baris Total Bobot, Skor, Total Skor (Format Baru) ---
    \multicolumn{2}{|l|}{\textbf{Total Bobot}} &                                                                                                                                         &                                       &             &            & 10                       \\
    \hline
    \multicolumn{2}{|l|}{\textbf{Skor}}        &                                                                                                                                         &                                       &             &            & 50                       \\
    \hline
    \multicolumn{2}{|l|}{\textbf{Total Skor}}  & \multicolumn{5}{c|}{50}                                                                                                                                                                                                               \\ % Mengganti "anjay" dengan "..."
\end{longtable}

Dari tabel \ref{tabel:uat_admin_cabang_hasil} di atas, kita dapat menghitung skor \textit{User Acceptance Test} (UAT) untuk perspektif Admin Cabang sebagai berikut:

\begin{align*}
    \text{Skor UAT Admin}(\%) & = \left( \frac{A}{B \times N} \right) \times 100\%   \\
                              & = \left( \frac{50}{50 \times 1} \right) \times 100\% \\
                              & = \left( \frac{50}{50} \right) \times 100\%          \\
                              & = 1 \times 100\%                                     \\
                              & = 100\%
\end{align*}

Kesimpulan dari hasil UAT perspektif Admin Cabang menunjukkan bahwa aplikasi Tamu.id mendapatkan skor 100\%, yang berarti aplikasi ini sangat diterima oleh Admin Cabang dan masuk dalam kategori Sangat Layak berdasar tabel \ref{tabel:kategori-kelayakan}.


\begin{longtable}{|p{0.8cm}|p{7.5cm}|>{\centering\arraybackslash}p{0.8cm}|>{\centering\arraybackslash}p{0.8cm}|>{\centering\arraybackslash}p{0.8cm}|>{\centering\arraybackslash}p{0.8cm}|>{\centering\arraybackslash}p{0.8cm}|}
    \caption{Hasil Kuesioner \textit{User Acceptance Test} - Perspektif \textit{SuperAdmin}} \label{tabel:uat_superadmin}                                                                                                                                                               \\
    \hline
    \multirow{2}{*}{\textbf{No.}}              & \multirow{2}{*}{\textbf{Pertanyaan}}                                                                                                     & \multicolumn{5}{c|}{\textbf{Jawaban}}                                                       \\ \cline{3-7}
                                               &                                                                                                                                          & \textbf{STS}                          & \textbf{TS} & \textbf{N} & \textbf{S} & \textbf{SS} \\
    \hline
    \endfirsthead

    \multicolumn{7}{c}%
    {{\bfseries Tabel \thetable\ lanjut dari halaman sebelumnya}}                                                                                                                                                                                                                       \\
    \hline
    \multirow{2}{*}{\textbf{No.}}              & \multirow{2}{*}{\textbf{Pertanyaan}}                                                                                                     & \multicolumn{5}{c|}{\textbf{Jawaban}}                                                       \\ \cline{3-7}
                                               &                                                                                                                                          & \textbf{STS}                          & \textbf{TS} & \textbf{N} & \textbf{S} & \textbf{SS} \\
    \hline
    \endhead

    \hline \multicolumn{7}{r}{\textit{Bersambung ke halaman berikutnya}}                                                                                                                                                                                                                \\
    \endfoot

    \hline % \endlastfoot
    \endlastfoot

    % --- Pertanyaan UAT untuk SuperAdmin (10 Pertanyaan) ---
    1.                                         & Apakah proses login ke dashboard SuperAdmin mudah dan aman?                                                                              &                                       &             &            &            & 1           \\
    \hline
    2.                                         & Apakah tampilan dashboard global SuperAdmin menyajikan informasi ringkasan sistem secara efektif dan mudah dipahami?                     &                                       &             &            &            & 1           \\
    \hline
    3.                                         & Apakah pengelolaan data master sistem (seperti daftar Cabang/Hotel, Kota, OTA global) mudah dilakukan?                                   &                                       &             &            &            & 1           \\
    \hline
    4.                                         & Apakah fitur untuk menambah, mengedit, dan mengelola akun staf (Admin Cabang, SuperAdmin lain) serta hak aksesnya intuitif?              &                                       &             &            &            & 1           \\
    \hline
    5.                                         & Apakah proses verifikasi identitas pengguna (KYC) dari sisi SuperAdmin berjalan dengan alur yang jelas dan efisien?                      &                                       &             &            & 1          &             \\
    \hline
    6.                                         & Apakah proses approval atau penolakan permintaan pencairan saldo referral pengguna mudah dikelola dan dilacak?                           &                                       &             &            & 1          &             \\
    \hline
    7.                                         & Apakah fitur untuk membuat dan mengelola voucher yang berlaku secara global atau untuk banyak cabang mudah digunakan?                    &                                       &             &            &            & 1           \\
    \hline
    8.                                         & Apakah sistem SuperAdmin menyediakan alat yang memadai untuk memantau aktivitas penting dan kesehatan sistem secara keseluruhan?         &                                       &             &            &            & 1           \\
    \hline
    9.                                         & Apakah laporan global atau data analitik yang dihasilkan sistem (jika ada) membantu dalam analisis dan pengambilan keputusan strategis?  &                                       &             &            &            & 1           \\
    \hline
    10.                                        & Secara keseluruhan, apakah antarmuka SuperAdmin efektif dalam membantu saya menjalankan tugas administrasi sistem dan pengawasan global? &                                       &             &            &            & 1           \\
    \hline
    % --- Baris Total Bobot, Skor, Total Skor (Format Baru) ---
    \multicolumn{2}{|l|}{\textbf{Total Bobot}} &                                                                                                                                          &                                       &             & 2          & 8                        \\
    \hline
    \multicolumn{2}{|l|}{\textbf{Skor}}        &                                                                                                                                          &                                       &             & 8          & 40                       \\
    \hline
    \multicolumn{2}{|l|}{\textbf{Total Skor}}  & \multicolumn{5}{c|}{48}                                                                                                                                                                                                                \\ % Mengganti "anjay" dengan "..."
\end{longtable}

Dari tabel \ref{tabel:uat_superadmin} di atas, kita dapat menghitung skor \textit{User Acceptance Test} (UAT) untuk perspektif SuperAdmin sebagai berikut:

\begin{align*}
    \text{Skor UAT SuperAdmin}(\%) & = \left( \frac{A}{B \times N} \right) \times 100\%   \\
                                   & = \left( \frac{48}{50 \times 1} \right) \times 100\% \\
                                   & = \left( \frac{48}{50} \right) \times 100\%          \\
                                   & = 0.96 \times 100\%                                  \\
                                   & = 96\%
\end{align*}

Kesimpulan dari hasil UAT perspektif SuperAdmin menunjukkan bahwa aplikasi Tamu.id mendapatkan skor 96\%, yang berarti aplikasi ini sangat diterima oleh SuperAdmin dan masuk dalam kategori Sangat Layak berdasar tabel \ref{tabel:kategori-kelayakan}.